% P16-C1: finite propositional abduction and explicit input boundaries.
\chapter{Finite Logical Abduction}
\label{sec:p16-finite-abduction}

Logical abduction searches for permitted assumptions which, together with a
background theory, consistently entail specified observations. This chapter
uses Eiter and Gottlob's finite propositional model: hypotheses and
observations are positive atoms, and an explanation is a subset of the
permitted hypotheses \cite[sec.~2, manuscript pp.~9--11]{EiterGottlobAbductionManuscript}.
The finite construction and examples below are supplied here with complete
proofs. No complexity classification is needed for this construction.

The source requires a consistent background theory. We also allow an
inconsistent theory as an input to the checking procedure, and prove that
this additional input case has no explanations. This extension does not
change the explanation criterion on the source's inputs.

\section{Background, hypotheses and observations}
\label{sec:p16-abduction-inputs}

\begin{hbdefinition}{def:p16-abduction-input}{Finite abduction input}
A \emph{finite abduction input} is a tuple
\(\mathcal P=(V,H,M,T)\), with the following data:
\begin{enumerate}
\item \(V\) is a finite subset of the propositional alphabet \(\mathsf{Prop}\).
\item \(H\subseteq V\) is the set of \emph{permitted hypotheses}, also
      called the \emph{abducible atoms}.
\item \(M\subseteq V\) is the set of \emph{observations}; Eiter and
      Gottlob call these \emph{manifestations}.
\item \(T\) is a finite set of classical propositional formulas whose
      occurring atoms all belong to \(V\), called the \emph{background
      theory}.
\end{enumerate}
The sets may be empty. Neither \(H\cap M=\varnothing\) nor
\(H\cup M=V\) is required. Hypotheses and observations in this definition
are atoms themselves: negative literals and compound formulas are not
members of \(H\) or \(M\). The theory \(T\) may contain negations
and compound formulas. Its formulas use the grammar and abbreviations of
\cref{sec:prop-syntax-semantics-01}.
The source's truth constant and biconditional are written here using
the book's abbreviations \(\topv\) and \(\leftrightarrow\).

A \emph{candidate} is a set \(E\subseteq H\). In \(T\cup E\),
the atoms in \(E\) are treated as formulas asserted positively.
An atom outside \(H\) cannot be inserted into a candidate merely because
it would imply an observation. Changing \(H\) changes the input problem.
\index{abduction!permitted hypotheses}\index{abduction!observations}
\end{hbdefinition}

\begin{hbdefinition}{def:p16-abduction-models}{Finite valuations and explanatory entailment}
Use the truth values \(\mathbb{B}=\{\mathbf{F},\mathbf{T}\}\) and the recursive
truth clauses of \cref{sec:prop-syntax-semantics-02}. Put
\[
 \mathsf{Val}(V)\defeq\{v:v\text{ is a function }V\to\mathbb{B}\}.
\]
For \(v\in\mathsf{Val}(V)\), evaluate a formula over \(V\) by those
same clauses, and write \(v\entails\varphi\) when its value is
\(\mathbf{T}\). For a set \(\Gamma\) of formulas over \(V\), put
\[
 \operatorname{Mod}_V(\Gamma)
   \defeq\{v\in\mathsf{Val}(V):
                  v\entails\gamma\text{ for every }\gamma\in\Gamma\}.
\]
We say \(\Gamma\) is \emph{consistent} in this chapter when this set
is nonempty, that is, when \(\Gamma\) is satisfiable. The notation
\(\Gamma\entails M\) abbreviates the assertion
\(\Gamma\entails m\) for every \(m\in M\), using the book's
classical semantic consequence relation. The next theorem identifies this
relation with the finite valuation test.

A candidate \(E\) is an \emph{explanation} for \(\mathcal P\) when
\[
 \operatorname{Mod}_V(T\cup E)\ne\varnothing
       \qquad\text{and}\qquad T\cup E\entails M.
\]
Write \(\operatorname{Expl}(\mathcal P)\) for the family of all such
explanations. This family corresponds to the source's
\(\operatorname{Sol}(\mathcal P)\).
\index{abduction!explanation}\index{abduction!consistency}
\end{hbdefinition}

\begin{hbtheorem}{thm:p16-abduction-semantics}{proved=annex:p16.abduction-semantics}{The finite criterion and its boundary cases}
For every finite abduction input and every \(E\subseteq H\),
\(E\) is an explanation exactly when
\(\operatorname{Mod}_V(T\cup E)\ne\varnothing\) and
\[
 v(m)=\mathbf{T}\qquad
  (v\in\operatorname{Mod}_V(T\cup E),\ m\in M).
\]
The use of functions on \(V\) agrees with the book's valuations on all
of \(\mathsf{Prop}\). After naming the finite alphabet in \(\mathsf{Prop}\) and
expanding those two abbreviations, inputs with consistent \(T\) have
precisely Eiter and Gottlob's explanation criterion
\cite[defs.~2.1--2.2, manuscript pp.~9--10]{EiterGottlobAbductionManuscript}.
The additional and empty cases satisfy:
\begin{enumerate}
\item If \(T\) is inconsistent, then
      \(\operatorname{Expl}(\mathcal P)=\varnothing\).
\item The empty set is an explanation exactly when \(T\) is consistent
      and \(T\entails M\).
\item If \(M=\varnothing\), the explanations are exactly those
      \(E\subseteq H\) for which \(T\cup E\) is consistent.
\end{enumerate}
\end{hbtheorem}

\begin{hbwarning}{def:p16-abduction-two-tests}{Consistency and sufficiency are separate tests}
\par\nobreak
Classical entailment from an inconsistent set is vacuous: there is no
satisfying valuation to falsify an observation. Such a set fails the
consistency test and gives no explanation. Conversely, consistency only
provides one model of the proposed assumptions. Explanation requires every
such model to satisfy every observation. A single countervaluation refutes
that requirement.

Membership in \(H\) is an additional input restriction. A consistent set
of atoms that entails \(M\) can still be ineligible because it includes
an atom outside \(H\). The examples below distinguish these failures.
\end{hbwarning}

\section{Two notions of a small explanation}
\label{sec:p16-abduction-minimality}

\begin{hbdefinition}{def:p16-abduction-minimality}{Subset-minimal and minimum-cardinality explanations}
An explanation \(E\) is \emph{subset-minimal} if no proper subset of
\(E\) is an explanation. It has \emph{minimum cardinality} if no
explanation has fewer members. Define the two families by
\[
\begin{aligned}
 \operatorname{MinExpl}(\mathcal P)
   &\defeq\{E\in\operatorname{Expl}(\mathcal P):\\
   &\hspace{1.5em}\text{no }F\subsetneq E\text{ is an explanation}\},\\
 \operatorname{CardExpl}(\mathcal P)
   &\defeq\{E\in\operatorname{Expl}(\mathcal P):\\
   &\hspace{1.5em}|E|\leq|F|\text{ for every explanation }F\}.
\end{aligned}
\]
If there are no explanations, both families are empty. These are the
subset and cardinality preference criteria of the source
\cite[def.~2.3 and following discussion, manuscript pp.~10--11]{EiterGottlobAbductionManuscript}.
They compare sets of assumptions. No probability, cost or causal relation
has been specified by either definition.
\index{abduction!subset-minimal explanation}
\index{abduction!minimum-cardinality explanation}
\end{hbdefinition}

\begin{hbtheorem}{thm:p16-abduction-minimality}{proved=annex:p16.abduction-minimality}{Existence, inclusion and the deletion criterion}
\par\nobreak
Every explanation contains a subset-minimal explanation. If any explanation
exists, both minimality families are nonempty, and
\[
 \operatorname{CardExpl}(\mathcal P)
        \subseteq\operatorname{MinExpl}(\mathcal P)
        \subseteq\operatorname{Expl}(\mathcal P).
\]
Distinct subset-minimal explanations are incomparable under inclusion;
their cardinalities may differ.

If \(E\) is an explanation, then
\[
\begin{aligned}
 E\in\operatorname{MinExpl}(\mathcal P)
 \quad\Longleftrightarrow\quad&
 E\setminus\{h\}\notin\operatorname{Expl}(\mathcal P)
                  \\
 &\text{for every }h\in E.
\end{aligned}
\]
This is the source's single-deletion characterization
\cite[prop.~2.1, p.~11]{EiterGottlobAbductionManuscript}.
It also holds for the extended input convention used here.

The relevant monotonicity has a consistency qualification: if \(E\) is
an explanation, \(E\subseteq F\subseteq H\), and \(T\cup F\) is
consistent, then \(F\) is an explanation. An arbitrary superset of an
explanation need not be an explanation, because it can become inconsistent.
\end{hbtheorem}

Subset-minimality can therefore be checked by testing each one-atom
deletion of a known explanation. Minimum cardinality compares it with all
explanations, including those incomparable with it. The unequal-size example
below makes the distinction explicit.

\section{An exhaustive finite construction}
\label{sec:p16-abduction-enumeration}

\begin{hbmethod}{def:p16-abduction-enumeration}{Enumerating explanations and both minimality families}
The input is the finite tuple \(\mathcal P=(V,H,M,T)\). List \(V\)
and \(H\) by increasing indices in \(\mathsf{Prop}=\{p_i:i\in\N\}\).
Put \(n=|V|\), \(r=|H|\), and initialize a list \(\mathcal L\)
with no members.
\begin{enumerate}
\item Enumerate every binary vector of length \(r\), in lexicographic
      order. Its candidate \(E\subseteq H\) contains the hypothesis
      at a position exactly when the bit at that position is one.
\item For this \(E\), initialize two Boolean flags to false:
      \emph{has a model} and \emph{has a countervaluation}. Enumerate
      every binary vector of length \(n\), in lexicographic order,
      interpreting zero as
      \(\mathbf{F}\) and one as \(\mathbf{T}\). This gives a valuation on
      \(V\). Evaluate every formula of \(T\) and every atom of \(E\).
      If all are true, set \emph{has a model} to true. If in addition some
      \(m\in M\) is false, set \emph{has a countervaluation} to true.
      Once set, a flag stays true for this candidate.
\item After the valuation scan, append \(E\) to \(\mathcal L\)
      exactly when \emph{has a model} is true and
      \emph{has a countervaluation} is false.
\item After all candidates, form \(\mathcal L_{\subseteq}\) by
      retaining exactly those members \(E\) of \(\mathcal L\) for
      which no other member is a proper subset of \(E\).
\item If \(\mathcal L\) is empty, put \(\mathcal L_{\#}\) empty.
      Otherwise compute the least cardinality of a member of \(\mathcal L\)
      by a finite scan, and let \(\mathcal L_{\#}\) retain exactly
      the members of that cardinality. Return all three lists.
\end{enumerate}
The empty binary vector occurs once when its length is zero. In particular,
empty \(H\) still gives the candidate \(\varnothing\), and empty
\(V\) still gives the unique empty valuation. Testing all observations
when \(M\) is empty cannot set the countervaluation flag.
\end{hbmethod}

\begin{hbtheorem}{thm:p16-abduction-enumeration}{proved=annex:p16.abduction-enumeration}{Termination and exact output}
The construction terminates. Viewed as sets of their listed members, its
outputs are exactly
\[
\begin{aligned}
 \mathcal L&=\operatorname{Expl}(\mathcal P),\\
 \mathcal L_{\subseteq}&=\operatorname{MinExpl}(\mathcal P),\\
 \mathcal L_{\#}&=\operatorname{CardExpl}(\mathcal P).
\end{aligned}
\]
No candidate or output member is duplicated. The unshortened nested scans
in steps 1--3 inspect exactly \(2^{r+n}\) candidate--valuation pairs.
This count does not make formula evaluation or the subsequent filtering
constant-time operations. No optimality or complexity-class claim is made.
\end{hbtheorem}

\section{Complete small examples}
\label{sec:p16-abduction-examples}

\begin{hbdefinition}{def:p16-abduction-examples}{Six explicit inputs}
In the following inputs, distinct letters among \(a,b,c,z,m\) denote
distinct members of \(\mathsf{Prop}\). The tuple order is always \((V,H,M,T)\).
These are additional handbook examples.
\[
\begin{aligned}
 \mathcal P_{\rm inconsistent}
   &\defeq(\{a,m\},\{a\},\{m\},\{a,\neg a\}),\\
 \mathcal P_{\rm given}
   &\defeq(\{a,m\},\{a\},\{m\},\{m\}),\\
 \mathcal P_{\rm empty}
   &\defeq(\varnothing,\varnothing,\varnothing,\varnothing),\\
 \mathcal P_{\rm sizes}
   &\defeq(\{a,b,c,m\},\{a,b,c\},\{m\},\\
   &\hspace{2em}\{a\to m,(b\land c)\to m\}),\\
 \mathcal P_{\rm constraint}
   &\defeq(\{a,b,m\},\{a,b\},\{m\},\\
   &\hspace{2em}\{a\to m,\neg(a\land b)\}),\\
 \mathcal P_{\rm permitted}
   &\defeq(\{a,z,m\},\{a\},\{m\},\{z\to m\}).
\end{aligned}
\]
\end{hbdefinition}

\begin{hbtheorem}{thm:p16-abduction-examples}{proved=annex:p16.abduction-examples}{All explanation families for the six inputs}
The complete families are as follows.
\begin{enumerate}
\item \emph{Inconsistent background.} All three explanation families of
      \(\mathcal P_{\rm inconsistent}\) are empty, even though every
      \(T\cup E\) entails \(m\) vacuously.
\item \emph{Observation already given.}
      \(\operatorname{Expl}(\mathcal P_{\rm given})
          =\{\varnothing,\{a\}\}\).
      Both minimality families equal \(\{\varnothing\}\).
\item \emph{Empty data.} All three families of
      \(\mathcal P_{\rm empty}\) equal \(\{\varnothing\}\).
      This is one explanation, rather than no explanations.
\item \emph{Unequal sizes.}
      \[
      \begin{aligned}
       \operatorname{Expl}(\mathcal P_{\rm sizes})
         &=\{\{a\},\{a,b\},\{a,c\},\{a,b,c\},\{b,c\}\},\\
       \operatorname{MinExpl}(\mathcal P_{\rm sizes})
         &=\{\{a\},\{b,c\}\},\\
       \operatorname{CardExpl}(\mathcal P_{\rm sizes})
         &=\{\{a\}\}.
      \end{aligned}
      \]
      The two subset-minimal explanations are incomparable, with sizes
      one and two.
\item \emph{Consistency constraint.} All three families of
      \(\mathcal P_{\rm constraint}\) equal \(\{\{a\}\}\).
      The candidate \(\{b\}\) is consistent with \(T\) but does
      not entail \(m\). The superset \(\{a,b\}\) of the explanation
      \(\{a\}\) fails consistency.
\item \emph{Permitted hypotheses.} All three families of
      \(\mathcal P_{\rm permitted}\) are empty. The set \(\{z\}\)
      is consistent with \(T\) and entails \(m\) with \(T\), but
      is not a candidate because \(z\notin H\). The permitted
      nonempty candidate \(\{a\}\) is consistent and non-explanatory.
\end{enumerate}
\end{hbtheorem}

\section{What an explanation establishes}
\label{sec:p16-abduction-interpretation}

\begin{hbtheorem}{thm:p16-abduction-truth-boundary}{proved=annex:p16.abduction-truth-boundary}{Explanatory assumptions need not follow from the observations}
In \(\mathcal P_{\rm sizes}\), both \(\{a\}\) and \(\{b,c\}\)
are subset-minimal explanations, but
\[
\begin{aligned}
 T\cup M&\not\entails a,\\
 T\cup M&\not\entails(b\land c).
\end{aligned}
\]
Indeed there is one valuation satisfying \(T\cup M\) which makes all
three permitted hypotheses false.
\end{hbtheorem}

An explanation asserts that, \emph{if} its hypotheses hold together with
the background, the observations follow, and that these combined assumptions
are satisfiable. The theorem shows why this does not establish that the
selected hypotheses are true. A declaration that \(M\) contains the
observations does not by itself supply evidence for their measurement or for
the background theory either.

\begin{hbwarning}{def:p16-abduction-causal-boundary}{Logical explanation and causal interpretation}
The arrows in \(T\) are classical material implications. The input
contains no causal relation, intervention semantics, temporal ordering or
probability distribution. Consequently the construction establishes a
logical property of the supplied formulas and permitted assumptions. A
claim about what actually caused an observation requires additional
structure and evidence. Choosing fewer assumptions, or removing redundant
ones, does not supply that structure.

This chapter supplies the finite atomic-hypothesis foundation and its stated
procedures. General formula hypotheses and their translations, first-order
abduction, other preference orders, the source's complexity classifications,
probabilistic abduction and causal explanation require separate treatments.
Consistency-based diagnosis, its complete conflict search, measurement
updates and default correspondence are supplied in
\cref{sec:p16-consistency-based-diagnosis}.
The exact distinction between consistency diagnosis and explanatory
entailment is proved in \cref{thm:p16-diag-abduction-boundaries}.
The distinction between generating an analogical candidate and
establishing its truth by entailment is
\cref{thm:p19-analogy-candidate-truth}; the associated task evidence
is supplied in \cref{sec:p19-analogy-and-problem-transfer}.
The separate deductive-explanation model, its verification
condition and its account of relative emergence are supplied in
\cref{sec:p16-deductive-nomological-explanation}.
The complete source
presentation and the broader explanation/abduction subject remain beyond
this bounded unit.
\end{hbwarning}
